\documentclass[czech,bachelor]{diploma}
\usepackage[autostyle=true,czech=quotes]{csquotes}
\usepackage[backend=biber, style=iso-numeric, alldates=iso]{biblatex}
\addbibresource{sources.bib}
\usepackage{dcolumn}
\usepackage{subfig}
\usepackage[cpp]{diplomalst}
\usepackage{chngcntr}
\usepackage{float}
\usepackage{array}
\usepackage{tikz}
\usepackage{listings}

\usetikzlibrary{positioning}
\lstdefinelanguage{JavaScript}{keywords={function, return, var, let, const, if, else, for, while, do, break, continue, switch, case, default, new, this, typeof, true, false, null, undefined}, keywordstyle=\bfseries, sensitive=true, comment=[l]{//}, morecomment=[s]{/*}{*/}, commentstyle=\itshape, morestring=[b]", morestring=[b]', stringstyle=\ttfamily}

\ThesisAuthor{Jakub Černík}
\ThesisSupervisor{Ing. Martin Kot, Ph.D.}
\CzechThesisTitle{Komponenta výukového serveru TI - Amortizovaná složitost 2}
\EnglishThesisTitle{Component of Teaching Server for Theoretical Computer Science - Amortized Complexity 2}

\SubmissionYear{2026}
\ThesisAssignmentFileName{ThesisSpecification_CER0548_vsboee25037FA1.pdf}

\Acknowledgement{PODEKOVANI.PODEKOVANIPODEKOVANIPODEKOVANIPODEKOVANIPODEKOVANIPODEKOVANIPODEKOVANIPODEKOVANIPODEKOVANIPODEKOVANIPODEKOVANI}

\CzechAbstract{Celkovým cílem této webové aplikace je zpřístupnit koncept amortizované složitosti prostřednictvím vizualizací a interaktivních simulací průběhu algoritmů, doplněných o podrobné vysvětlení jednotlivých kroků. Aplikace zároveň obsahuje rozsáhlé teoretické pasáže, které přispívají k hlubšímu porozumění dané problematiky. Na webu jsou prezentovány tři algoritmy běžně probírané ve výuce informatiky. Kombinací těchto prvků vzniká nástroj, který studentům poskytuje srozumitelné a systematické vysvětlení daného učiva.}

\CzechKeywords{DOPLNIT}
\EnglishAbstract{DOPLNIT}
\EnglishKeywords{DOPLNIT}

\newcolumntype{d}[1]{D{,}{,}{#1}}
\newcolumntype{L}[1]{>{\raggedright\arraybackslash}p{#1}}
\counterwithin{figure}{section}
\counterwithin{table}{section}
\setcounter{tocdepth}{2}
\raggedbottom
\begin{document}

\MakeTitlePages

\chapter{Úvod}

V této kapitole uveďte motivaci práce, hlavní cíle a stručně popište strukturu celého dokumentu.

\section{Motivace a kontext práce}
Stručně popište, proč je amortizovaná složitost pro studenty obtížná a proč je vhodné ji vysvětlovat pomocí interaktivních vizualizací.

\section{Cíle bakalářské práce}
Uveďte hlavní cíl práce a dílčí cíle teoretické i praktické části. Doplňte, jaké konkrétní výstupy má práce přinést.

\section{Metodika a postup řešení}
Vysvětlete, jak jste postupoval od teoretické analýzy přes návrh komponenty až po implementaci a ověření funkčnosti.

\section{Struktura práce}
Krátce popište obsah jednotlivých kapitol a jejich návaznost.

\chapter{Základní problematika}

Tato kapitola bude doplněna.

\chapter{Návrh komponenty výukového serveru}

Tato kapitola bude doplněna.

\chapter{Implementace komponenty}

Tato kapitola bude doplněna.

\chapter{Závěr}

V závěru zhodnoťte dosažené výsledky práce a uveďte jejich přínos.

\section{Vyhodnocení splnění cílů}
Po jednotlivých bodech uveďte, které cíle práce byly splněny a jakým způsobem.

\section{Hlavní přínosy práce}
Stručně shrňte přínos pro výuku amortizované složitosti a praktické využití vytvořené komponenty.

\section{Omezení řešení}
Uveďte, co je v aktuální verzi zjednodušeno nebo mimo rozsah práce.

\section{Možnosti budoucího rozšíření}
Navrhněte smysluplné směry dalšího vývoje komponenty a navazujících prací.

\end{document}